\chapter{INTRODUÇÃO}
\label{introducao}

\section{DELIMITAÇÃO DO PROBLEMA}

Os parágrafos não contêm recuo. O espaçamento entre linhas é de 1,5 linhas e entre os parágrafos deve ser mantido o espaçamento de 18 pt.

\section{HIPÓTESES}

Descreva aqui as hipóteses da pesquisa.

\section{OBJETIVOS}

Descreva aqui os objetivos da pesquisa.

\section{JUSTIFICATIVA}

Descreva aqui a justificativa da pesquisa.

\begin{quadro}[htbp]
  \caption{Componentes relacionados à segurança viária, intervenções e benefícios}
  \label{quad:componentes}
  \centering
  \small
  \begin{tabular}{|p{3.5cm}|p{5cm}|p{5cm}|}
    \hline
    \textbf{COMPONENTE} & \textbf{INTERVENÇÕES} & \textbf{BENEFÍCIOS} \\
    \hline
    Controle de Velocidade & Estabelecer e impor leis de limite de velocidade em todo o país. & Redução das lesões e mortes no trânsito, além dos custos socioeconômicos relacionados. \\
    \hline
    Controle de Velocidade & Construir ou modificar vias para acalmar o tráfego, como rotatórias, estreitamentos ou blitz. & Melhoria na poluição do ar, no consumo de combustível e na poluição sonora. \\
    \hline
    Controle de Velocidade & Exigir que os fabricantes de automóveis instalem novas tecnologias que ajudem a lembrar os motoristas dos limites de velocidade. & Ambiente mais amigável para práticas de caminhada e ciclismo. \\
    \hline
  \end{tabular}
  \\[0.2cm]
  {\small Fonte: Camargo e Souza, 2018.}
\end{quadro}

\begin{figure}[htbp]
  \centering
  % \includegraphics[width=0.8\textwidth]{imagens/desenho_metodologico.png}
  \caption{Desenho metodológico da pesquisa}
  \label{fig:desenho_metodologico}
  \vspace{0.2cm}
  \par
  {\small Fonte: Camargo e Souza, 2018.}
\end{figure}

\section{ESTRUTURA DA DISSERTAÇÃO}

O trabalho está organizado como segue. O Capítulo \ref{capitulo1} apresenta detalhes sobre a fundamentação e segurança viária. O Capítulo \ref{capitulo2} apresenta os materiais e métodos utilizados. O Capítulo \ref{capitulo3} descreve a aplicação e análise dos resultados. Por fim, o Capítulo \ref{conclusao} apresenta as conclusões e trabalhos futuros propostos.